\section{Positive filters and intervals at every pin}\label{sec:filters}
The arguments in this section apply whenever
\begin{equation}\label{eq:S}
 S=\sum_{j>j_0}2^{n_{j-1}-aL_j}<\infty.
\end{equation}
This follows from superlacunarity, since eventually $aL_j\geq2n_{j-1}$
and $n_{j-1}\geq j-1$. For a fixed ratio it follows from $aK>1$, since
the summands are $2^{-(aK-1)n_{j-1}}$.

\begin{proposition}[One fixed positive restriction]\label{prop:fixed-filter}
For every $p\in P$ and fixed $0<\varepsilon\leq1$, the actual measure
$\nu_{p,\varepsilon}=(D_p)_*(M_\varepsilon(p,\cdot)\mu)$ has a
jointly continuous density $g_\varepsilon(p,r)$, with
\begin{equation}\label{eq:fixed-good-bound}
 \norm{g_\varepsilon(p,\cdot)}_\infty\leq C_*/\varepsilon,
 \qquad \nu_{p,\varepsilon}(\R)\geq1-K_X\varepsilon^\beta.
\end{equation}
One possible constant is
\begin{equation}\label{eq:Cstar}
 C_*=(D_0/v_0)2^{n_{j_0}}+4D_0(1+v_0^{-1})S.
\end{equation}
It may depend on the finitely many initial stages, but is independent of
$p$ and any finest observation scale.
\end{proposition}
\begin{proof}
Let $g_{j,p,\varepsilon}$ be the masked density for $\mu_j$.
Apply Lemma~\ref{lem:fixed-band} in every parent square, multiply by
its exact parent mass, and add. The masses sum to one. By $Bw=A$,
\begin{equation}\label{eq:fixed-increment}
 \norm{g_{j,p,\varepsilon}-g_{j-1,p,\varepsilon}}_\infty
 \leq4D_0(v_0^{-1}+\varepsilon^{-1})2^{n_{j-1}-aL_j}.
\end{equation}
No count of cells appears. At the initial filled stage, vertical coarea
on each square gives
$\norm{g_{j_0,p,\varepsilon}}_\infty\leq(D_0/v_0)2^{n_{j_0}}$.
Thus the densities converge in $L^\infty$ uniformly in $p$, and the
bound~\eqref{eq:Cstar} follows.

For a fixed $p$, the discontinuities of the mask are on
$y_1=p_1\pm\varepsilon$. These lines have zero $\mu$-mass by
\eqref{eq:horizontal-frostman}. Weak convergence can therefore be applied
to a bounded continuous distance test function multiplied by the mask.
It identifies the weak limit with $(D_p)_*(M_\varepsilon\mu)$.
The common bounded output support makes the $L^\infty$ convergence also
$L^1$ convergence, identifying its density with the same measure.
This argument works for \emph{each} pin with the same bounds, without an
exceptional pin set. The retained-mass estimate follows directly from
\eqref{eq:horizontal-frostman}.

To obtain the continuous version, write the finite filled-tail density
$f_j$ as a finite linear combination of rectangle indicators. Polar coarea
gives, for $r>0$,
\begin{equation}\label{eq:polar-continuity}
 g_{j,\varepsilon}(p,r)=r\int_0^{2\pi}
 f_j(p+r\omega(\theta))
 \ind_{\{|r\cos\theta|\geq\varepsilon\}}\,d\theta,
 \quad\omega(\theta)=(\cos\theta,\sin\theta).
\end{equation}
A circle with positive radius meets each rectangle boundary in only finitely
many angles. The mask equality also has at most finitely many angles.
For a convergent sequence of $(p,r)$ the integrand therefore converges at
almost every angle and has a bounded integrable majorant. Dominated
convergence proves joint continuity. Around $r=0$ the density is zero by
uniform vertical separation; extend by zero for negative $r$.

For continuous radius fibers the supremum is the essential supremum.
Thus~\eqref{eq:fixed-increment} is a uniform supremum estimate on the
continuous versions over $P\times\R$. Their uniform limit is jointly
continuous and represents the actual positive restriction at every pin.
\end{proof}

\begin{corollary}[Raw absolute continuity and weak integrability]
\label{cor:weak}
Every $\nu_p$ is absolutely continuous. Its density satisfies
\begin{equation}\label{eq:weak-alpha}
 \sup_{p\in P}|\{r:g_p(r)>T\}|\leq C T^{-\alpha}
 \qquad(T>0).
\end{equation}
In particular $\sup_p\norm{g_p}_q<\infty$ for $1\leq q<\alpha$.
\end{corollary}
\begin{proof}
As $\varepsilon\downarrow0$, the positive restricted measures increase
to $\nu_p$, since $\mu\{y_1=p_1\}=0$. A Lebesgue-null set has zero
mass for every restriction and hence for the limit. This proves absolute
continuity. Write $g_p=g_\varepsilon(p,\cdot)+b_{p,\varepsilon}$ with
$b_{p,\varepsilon}\geq0$ and
$\int b_{p,\varepsilon}\leq K_X\varepsilon^\beta$.
For sufficiently large $T$, choose $\varepsilon=2C_*/T\leq1$.
Then the good density is at most $T/2$, so Markov's inequality gives
\[
 |\{g_p>T\}|\leq\frac2T\int b_{p,\varepsilon}
 \leq2K_X(2C_*)^\beta T^{-(1+\beta)}.
\]
Since $1+\beta=\alpha$, this is the claim for large $T$; bounded output
support handles smaller $T$. Layer-cake integration gives the stated
strong $L^q$ bounds.
\end{proof}

\begin{proposition}[Uniform length in the actual distance sets]
\label{prop:intervals}
There is $\ell>0$ such that every $D_p(E_+)$, $p\in P_-$, contains
an interval of length $\ell$. The reversed conclusion holds for
$p\in P_+$ and source $E_-$.
\end{proposition}
\begin{proof}
Choose a fixed $\varepsilon_0>0$ with
$K_X\varepsilon_0^\beta\leq1/2$. Every nonnegative fiber of
$g_{\varepsilon_0}$ then has mass at least $1/2$ on the interval
$I=[v_0,D_0]$ of length $W=D_0-v_0$. It has some point $r_p$ with
$g_{\varepsilon_0}(p,r_p)\geq1/(2W)$. Joint continuity and compactness,
using the continuous zero extension outside $I$, supply one $\delta>0$
such that
\[
 |r-r_p|\leq\delta\quad\Longrightarrow\quad
 g_{\varepsilon_0}(p,r)\geq1/(4W)
\]
for every $p$. The support of this submeasure is contained in the compact
set $D_p(E_+)$ because the submeasure is a positive restriction of the
actual source law. Thus the entire interval is in the actual image, not
only in a finite approximation or an almost-everywhere version of that
image. Taking any fixed $\ell\leq\delta$ proves the first assertion.

In the reverse configuration the vertical separation has negative sign,
but the absolute derivative bound is unchanged. The vertical branch is
still unique and the BV monotonicities simply reverse sign. The same proof
therefore applies. If necessary take the smaller of the two length bounds.
\end{proof}
Proposition~\ref{prop:intervals} proves Theorem~\ref{thm:main-intervals}.
In particular, failure of $L^2$ or even unboundedness of the raw density
is compatible with an interval in every one of the distance images here.
